\section*{Capítulo \ref{ch:b3:computational-chemistry} --- Química computacional: Hartree--Fock e DFT}

\begin{solution}{exo:b3:computational-chemistry:1}
$\dd E/\dd\alpha = \frac32 - \sqrt{2/\pi\alpha} = 0$ dá $\alpha = 8/9\pi = 0.283\,
a_0^{-2}$ e $E = \frac32\alpha - 2\sqrt{2\alpha/\pi} = -4/3\pi = \qty{-0.4244}{\hartree}$,
15\,\% acima do valor exato $-0.5$: uma única gaussiana não tem nem o bico nem
a cauda da função $1s$.
\end{solution}

\begin{solution}{exo:b3:computational-chemistry:2}
STO-3G: 5 funções por C ($1s$, $2s$, três $2p$) e 1 por H: $6 \times 5 + 6 =
36$. 6-31G(d): por C, 1 (caroço) + $2 \times 4$ (valência dividida) + 6 ($d$) $= 15$;
por H, 2: $6 \times 15 + 6 \times 2 = 102$.
\end{solution}

\begin{solution}{exo:b3:computational-chemistry:3}
O hamiltoniano eletrônico não contém as massas nucleares: dentro da
\omterm{def:b3:computational-chemistry:born-oppenheimer}{aproximação de Born--Oppenheimer}, a SEP, logo $r_e$ e $k = U''(r_e)$, é a
mesma. $\tilde\omega = \sqrt{k/\mu}/2\pi c$ depende da \omterm{def:b3:quantum-model-systems:oscillator}{massa reduzida}, que
dobra: razão $\sqrt2 = 1.414$ (medida: 1.413).
\end{solution}

\begin{solution}{exo:b3:computational-chemistry:4}
Oito átomos dão $3N - 6 = 18$ vibrações: 17 reais e uma imaginária. Um
\omterm{def:b3:quantum-model-systems:operator}{autovalor} negativo da hessiana: um ponto de sela de primeira ordem, uma
estrutura de transição (aqui, por exemplo, de uma rotação ou de uma eliminação).
\end{solution}

\begin{solution}{exo:b3:computational-chemistry:5}
$\zeta = 27/16$, $E = -(27/16)^2 = \qty{-2.8477}{\hartree} = \qty{-77.49}{eV}$.
Erro $-2.8477 + 2.9034 = \qty{0.0557}{\hartree} = \qty{1.52}{eV}$, ou seja,
1.9\,\% da energia total --- mas maior que muitas energias de reação.
\end{solution}

\begin{solution}{exo:b3:computational-chemistry:6}
$(-1 - E)(-0.5 - E) - (-0.2 - 0.3E)^2 = 0$, isto é, $0.91E^2 + 1.38E + 0.46 = 0$:
$E = \qty{-1.0217}{\hartree}$ e \qty{-0.4947}{\hartree}. Sim: misturar a
segunda função abaixa a energia para baixo de $H_{11}$, como o \omterm{thm:b3:computational-chemistry:variational}{princípio
variacional} permite.
\end{solution}

\begin{solution}{exo:b3:computational-chemistry:7}
$(102/36)^4 = 64$: cerca de 640 minutos, umas 11 horas.
\end{solution}

\begin{solution}{exo:b3:computational-chemistry:8}
$0.5782 \times 27.211 = \qty{15.73}{eV}$, 0.30 eV acima do valor medido. Orbitais
congelados: o íon real relaxa e fica mais baixo, logo o valor de Koopmans é alto
demais. Correlação ausente: a molécula neutra (dois elétrons) tem mais \omterm{def:b3:computational-chemistry:correlation}{energia de
correlação} que o íon (um elétron), o que torna o valor verdadeiro maior.
Aqui o primeiro erro prevalece.
\end{solution}

\begin{solution}{exo:b3:computational-chemistry:9}
$-0.5960 - 2(-0.4666) = \qty{0.3372}{\hartree} = \qty{9.18}{eV}$ acima de dois átomos.
A função RHF mantém 50\,\% de \ce{H+ H-}, e separar um próton e um hidreto
custa a diferença entre a energia de ionização de H e a afinidade eletrônica
de H, que entra na média da energia.
\end{solution}

\begin{solution}{exo:b3:computational-chemistry:10}
$E(c_1^2 + 2c_1c_2S + c_2^2) = c_1^2H_{11} + 2c_1c_2H_{12} + c_2^2H_{22}$.
Derivando em relação a $c_1$ com $\partial E/\partial c_1 = 0$:
$E(2c_1 + 2c_2S) = 2c_1H_{11} + 2c_2H_{12}$, isto é, $(H_{11} - E)c_1 + (H_{12} -
ES)c_2 = 0$; do mesmo modo, $(H_{12} - ES)c_1 + (H_{22} - E)c_2 = 0$.
\end{solution}

\begin{solution}{exo:b3:computational-chemistry:11}
$k \approx [E(1.30) - 2E(1.35) + E(1.40)]/(0.05)^2 = 0.001417/0.0025 =
\qty{0.567}{\hartree\per\bohr\squared} = \qty{882}{N/m}$. Com $\mu = m_H/2 =
\qty{8.37e-28}{kg}$, $\omega = \sqrt{k/\mu} = \qty{1.03e15}{s^{-1}}$ e
$\tilde\omega = \omega/2\pi c = \qty{5452}{cm^{-1}}$, 24\,\% acima do
$\tilde\omega_e$ medido: uma curva RHF em \omterm{def:b3:computational-chemistry:basis}{base mínima} é íngreme demais (sem
correlação, base pequena demais), e é por isso que as frequências calculadas
são reescaladas para baixo.
\end{solution}

\begin{solution}{exo:b3:computational-chemistry:12}
Diferenças de alguns \unit{kJ/mol} que envolvem uma ligação de hidrogênio exigem
correlação e dispersão: um funcional híbrido com correção de dispersão, ou,
melhor, um método correlacionado de função de onda para as energias finais,
com uma base polarizada de valência tripla que inclua \omterm{def:b3:computational-chemistry:split-valence}{funções difusas}. Otimize
os dois confôrmeros no mesmo nível, confirme os mínimos pelas frequências
(todas reais), acrescente as \omterm{def:b3:quantum-model-systems:oscillator}{energias do ponto zero}, teste a base com um
cálculo maior e compare com um sistema aparentado de energia conformacional conhecida.
\end{solution}

\begin{solution}{pb:b3:computational-chemistry:1}
\textbf{1.} Cada \omterm{def:b3:computational-chemistry:basis}{orbital do tipo Slater} de uma \omterm{def:b3:computational-chemistry:basis}{base mínima} é substituído por uma
contração fixa de três gaussianas ajustadas a ele.
\textbf{2.} O núcleo do hélio atrai mais seus elétrons, logo sua função $1s$
é mais compacta (expoentes maiores). $6.3624/3.4253 = 1.857 = (1.69/1.24)^2$.
\textbf{3.} Duas funções de base, dois elétrons, um orbital ocupado (e um
virtual).
\textbf{4.} $V_{\mathrm{nn}} = Z_{\mathrm{He}}Z_{\mathrm H}/R = 2/1.4632 =
\qty{1.3669}{\hartree}$.
\textbf{5.} As duas funções se sobrepõem fortemente (54\,\%): os átomos estão
próximos o bastante para se ligar.
\textbf{6.} $h_{11}$ é a energia de um elétron na função do He com os dois
núcleos, mas sem nenhum outro elétron; o núcleo de He (carga 2) o retém com
muito mais força.
\textbf{7.} $(h_{11} - \varepsilon)(h_{22} - \varepsilon) - (h_{12} - \varepsilon S)^2 = 0$:
$0.71185\varepsilon^2 + 2.79292\varepsilon + 2.44969 = 0$.
\textbf{8.} $\varepsilon = \qty{-2.600}{\hartree}$ e \qty{-1.324}{\hartree}.
\textbf{9.} $\mathbf h$ ignora a repulsão entre os dois elétrons; o orbital de
caroço fica contraído demais e baixo demais.
\textbf{10.} $F_{\mu\nu} = h_{\mu\nu} + \sum_{\lambda\sigma}P_{\lambda\sigma}[(\mu\nu|\sigma\lambda) -
\frac12(\mu\lambda|\sigma\nu)]$, com $P_{\lambda\sigma} = 2C_{\lambda1}C_{\sigma1}$.
\textbf{11.} A repulsão coulombiana de cada elétron pela nuvem de carga do
outro (e, em geral, a troca).
\textbf{12.} $\mathbf F$ depende de $\mathbf P$, que depende dos orbitais obtidos
a partir de $\mathbf F$: as equações são não lineares e são resolvidas por
aproximações sucessivas até a autoconsistência.
\textbf{13.} $c_1^2 + c_2^2 + 2c_1c_2S = 0.7684 + 0.0410 + 0.1906 = 1.0000$.
\textbf{14.} He: $1.5369 + 0.1906 = 1.727$; H: $0.0820 + 0.1906 = 0.273$
elétron.
\textbf{15.} Cargas: He $+0.27$, H $+0.73$. A carga positiva fica sobretudo no
hidrogênio: \ce{HeH+} é um próton ligado a um átomo de hélio, \ce{He + H+}, como
esperado, já que a energia de ionização de He (24.6 eV) supera de longe a de H (13.6 eV).
\textbf{16.} $-\varepsilon_1 = \qty{1.633}{\hartree} = \qty{44.4}{eV}$.
\textbf{17.} $E_{\mathrm{el}} = -2.8418 - 1.3669 = \qty{-4.2087}{\hartree}$.
\textbf{18.} Três ciclos chegam a \qty{-2.8418}{\hartree}. Cada ciclo fornece um
determinante, cuja energia é um limite superior da energia Hartree--Fock
convergida (\omterm{thm:b3:computational-chemistry:variational}{princípio variacional}), e as iterações a melhoram.
\textbf{19.} A base com $\zeta = 2.0925$ dá a energia mais baixa, logo é a
melhor para essa molécula (\omterm{thm:b3:computational-chemistry:variational}{princípio variacional}): a função do hélio é mais
compacta em \ce{HeH+} que no átomo livre, cujo $\zeta = 1.69$ a base padrão
copia.
\textbf{20.} O orbital antiligante vazio $\sigma^*$; sua energia aproximaria o
oposto da afinidade eletrônica de \ce{HeH+} (mal, em uma base tão pequena).
\textbf{21.} Zero (nenhum elétron); $\qty{-2.8078}{\hartree}$.
\textbf{22.} $-2.8078 - (-2.8418) = \qty{0.0340}{\hartree} = \qty{89}{kJ/mol}$.
\textbf{23.} Cerca da metade do valor medido: uma \omterm{def:b3:computational-chemistry:basis}{base mínima} não consegue
polarizar a densidade do hélio em direção ao próton (não há funções $p$), de
modo que a ligação fica fraca demais. (As correções de ponto zero e a \qty{298}{K}
são pequenas diante disso.)
\textbf{24.} $-(24.5874 + 54.4178)/27.211 = \qty{-2.9034}{\hartree}$; o átomo STO-3G
fica \qty{0.0956}{\hartree} (\qty{2.6}{eV}) alto demais.
\textbf{25.} As geometrias dependem de como a energia varia com $R$; a maior
parte do erro, concentrada no caroço, é quase a mesma em todas as geometrias e
se cancela.
\textbf{26.} \textbf{$E(\ce{HeH+}) = \qty{-2.8418}{\hartree}$ em RHF/STO-3G (base
padrão), a \qty{1.4632}{\bohr}} (e \qty{-2.8607}{\hartree} com o clássico
$\zeta_{\mathrm{He}} = 2.0925$).
\end{solution}
